\begin{helper}
Fix one pair of support labels \(B,J\) in a fixed history cell, and write
\(U\) for the terminal coordinate set with nodup reveal order \(o\).  Let
\(\mathcal P\) be the transcript family.  For each realized
branch/transcript pair \((\beta,\tau)\), write
\[
P_\tau,\qquad A_\tau,\qquad Z_\tau,\qquad F_{\beta,\tau}
\]
for the forced-present set, forced-absent set, selected absent set, and the
fixed present support \(B\cup R_{\beta,J}\).  Assume the realized geometry
\[
F_{\beta,\tau}\subseteq P_\tau,\qquad
P_\tau,A_\tau\subseteq U,\qquad
P_\tau\cap A_\tau=\emptyset,\qquad
Z_\tau\subseteq A_\tau ,
\]
the forward compatibility consequences for complete terminal assignments,
and functional disjointness for complete assignments compatible with two
realized cells.

Assume moreover the fixed residual scan interface supplied by
\(\noderef{FixedSetHistoryCellFixedBJResidualScanInterface}\): there are
one residual scan \(\operatorname{scan}_{B,J}\) and one reconstructed full
scan \(\operatorname{scan}^{\rm full}_{B,J}\), with the
delete-and-reinsert equality, reconstructed scan stability,
scan-to-compatibility, compatibility-to-scan, and residual functionality
clauses.

Then the two scan facts needed by the fixed-count red summation hold.
First, every residual assignment in the cylinder
\[
P_\tau\setminus F_{\beta,\tau}\subseteq A_{\rm res},\qquad
A_{\rm res}\cap(A_\tau\setminus Z_\tau)=\emptyset
\]
reconstructs to a complete assignment compatible with \((\beta,\tau)\):
\[
\operatorname{Comp}\bigl(
  (A_{\rm res}\setminus Z_\tau)\cup F_{\beta,\tau},
  \beta,\tau\bigr).
\]
Second, the same residual assignment has at most one realized owner: if
\[
\operatorname{Comp}\bigl(
  (A_{\rm res}\setminus Z_\tau)\cup F_{\beta,\tau},
  \beta,\tau\bigr)
\]
and
\[
\operatorname{Comp}\bigl(
  (A_{\rm res}\setminus Z_{\tau'})\cup F_{\beta',\tau'},
  \beta',\tau'\bigr)
\]
hold for realized pairs \((\beta,\tau)\) and \((\beta',\tau')\), then
\(\beta=\beta'\) and \(\tau=\tau'\).
\end{helper}
\begin{proof}
Unpack the fixed scan interface from
\(\noderef{FixedSetHistoryCellFixedBJResidualScanInterface}\).  It gives
the residual scan \(\operatorname{scan}_{B,J}\), the reconstructed full
scan \(\operatorname{scan}^{\rm full}_{B,J}\), and the five scan clauses
listed in the statement.

Let \(A_{\rm res}\subseteq U\) lie in the residual cylinder of a realized
\((\beta,\tau)\), and put
\[
A^\ast=(A_{\rm res}\setminus Z_\tau)\cup F_{\beta,\tau}.
\]
The realized geometry gives \(F_{\beta,\tau}\subseteq P_\tau\).  By the
delete-and-reinsert clause, the residual scan at \(A_{\rm res}\) is the
reconstructed full scan at
\[
(A_{\rm res}\setminus Z_\tau)\cup F_{\beta,\tau}.
\]
The reconstructed-stability clause evaluates that reconstructed full scan
on the same complete assignment.  Combining the two equalities gives
\[
\operatorname{scan}_{B,J}(A_{\rm res})
=(\beta,\tau).
\]
Applying the scan-to-compatibility clause to this output gives
\[
\operatorname{Comp}(A^\ast,\beta,\tau),
\]
which is exactly residual-cylinder completeness.

For single-owner uniqueness, suppose the same residual assignment
\(A_{\rm res}\) reconstructs to compatible complete assignments for two
realized pairs \((\beta,\tau)\) and \((\beta',\tau')\).  Apply the
compatibility-to-scan clause to the first compatibility certificate:
\[
\operatorname{scan}_{B,J}(A_{\rm res})=(\beta,\tau).
\]
Apply it again to the second compatibility certificate:
\[
\operatorname{scan}_{B,J}(A_{\rm res})=(\beta',\tau').
\]
The residual functionality clause of the same fixed scan applies to these
two outputs for the same residual assignment.  Hence
\(\beta=\beta'\) and \(\tau=\tau'\).  This proves the two residual facts
used directly by the released-block construction.
\end{proof}
